% 03-terms.tex — Clause 3: Terms, definitions and symbols
\chapter{Terms, definitions, and symbols}
\label{sec:terms}
\index{terms and definitions}

\pnum
For the purposes of this document, the following terms and definitions
apply.  Terms defined here are used with these precise meanings throughout
all normative clauses.  Where a term is used in its ordinary English sense
rather than as a defined term, it appears without emphasis.

\begin{note}
  Cross-references of the form \textit{(see \S\,N)} point to the clause
  where the concept is treated normatively.
\end{note}

% -------------------------------------------------------------------------
\section{General language terms}
\label{sec:terms.general}
% -------------------------------------------------------------------------

\subsection*{abstract syntax tree (AST)}
The in-memory tree representation of a Lain program produced by the parser
and consumed by semantic analysis.

\subsection*{behavior}
External appearance or action of a program as observed by a user or
another program.  See also: \textit{defined behavior},
\textit{implementation-defined behavior}, \textit{unspecified behavior},
\textit{undefined behavior}.

\subsection*{conforming implementation}
An implementation that correctly processes all programs that satisfy the
rules of this specification and issues a diagnostic message for every
program that violates a \textbf{Constraint} (see \S\,\ref{sec:conformance}).

\subsection*{conforming program}
A program that uses only the syntax and semantics defined in this
specification, does not rely on undefined behavior, and does not depend on
implementation-defined extensions.

\subsection*{constraint}
A requirement imposed on the source text.  If a program violates a
constraint, the implementation \textit{shall} emit a diagnostic message.
Constraints appear in \textbf{Constraint} environments throughout this
specification.

\subsection*{defined behavior}
Behavior that is fully and precisely specified by this document for every
conforming program.  Contrast with \textit{undefined behavior} and
\textit{implementation-defined behavior}.

\subsection*{diagnostic message}
\index{diagnostic message}
A human-readable message issued by the implementation to report a
constraint violation.  Every diagnostic shall include a source location
(file name, line number) and an error code of the form \texttt{[EXXX]} or
\texttt{[VRA]} (see Annex~B).

\subsection*{ill-formed program}
A program that violates one or more of the syntactic or semantic rules of
this specification.  An ill-formed program is not a conforming program and
a conforming implementation shall reject it with a diagnostic.

\subsection*{implementation}
The system (compiler, linker, runtime) that translates a Lain source program
into an executable.  The reference implementation is the \textit{lain}
compiler.

\subsection*{implementation-defined behavior}
\index{implementation-defined behavior}
Behavior that is valid but not fully specified by this document; the
implementation shall determine the exact behavior and document it.

\subsection*{program}
A collection of one or more Lain source files (\texttt{.ln}) that together
constitute a translation unit or a set of modules.

\subsection*{source file}
\index{source file}
A UTF-8 encoded text file with the extension \texttt{.ln} containing Lain
source text.

\subsection*{translation output}
The C99 source code emitted by a conforming implementation for a
well-formed Lain program.

\subsection*{undefined behavior}
\index{undefined behavior}
Behavior for which this specification imposes no requirements.  A program
that exhibits undefined behavior is erroneous.  In the safe fragment of
Lain, no operation produces undefined behavior.  Undefined behavior may
arise only inside \lain{unsafe} blocks or through direct C
interoperability.

\subsection*{unspecified behavior}
\index{unspecified behavior}
Behavior for which this specification provides two or more possibilities
without prescribing which shall occur.  A conforming program shall not
depend on unspecified behavior.

\subsection*{well-formed program}
A program that satisfies all syntactic and semantic rules of this
specification and is therefore accepted by a conforming implementation
without diagnostic.

% -------------------------------------------------------------------------
\section{Type system terms}
\label{sec:terms.types}
% -------------------------------------------------------------------------

\subsection*{algebraic data type (ADT)}
\index{algebraic data type}
A composite type that is either a product type (\textit{struct}) or a sum
type (\textit{enum}).

\subsection*{constrained type}
\index{constrained type}
A type annotated with a value range constraint of the form
\textit{T op bound}, e.g.\ \lain{int >= 0}
(see \S\,\ref{sec:07-types}).

\subsection*{integer rank}
\index{integer rank}
A total ordering on integer types that determines the result type of
mixed-type arithmetic.  The rank system is defined in
\S\,\ref{sec:07-types}.

\subsection*{integer widening}
\index{integer widening}
An implicit conversion from an integer type of lower rank to one of
higher rank, permitted in arithmetic and assignment contexts
(see \S\,\ref{sec:07-types}).

\subsection*{linear type}
\index{linear type}
A type whose values must be used exactly once.  Struct types and resource
types are linear.  Primitive types (\lain{int}, \lain{bool}, etc.) are
unrestricted (see \S\,\ref{sec:11-ownership}).

\subsection*{nominal type}
A type identified by its declared name rather than its structure.

\subsection*{ownership mode}
\index{ownership mode}
An annotation on a variable or parameter that specifies the kind of
access the holder has: \lain{shared} (read-only), \lain{var}
(read-write), or owned (no annotation on a linear type, transferred with
\lain{mov}) (see \S\,\ref{sec:11-ownership}).

\subsection*{primitive type}
\index{primitive type}
One of the built-in types: \lain{int}, \lain{u8}, \lain{u16}, \lain{u32},
\lain{u64}, \lain{usize}, \lain{i8}, \lain{i16}, \lain{i32}, \lain{i64},
\lain{isize}, \lain{f32}, \lain{f64}, \lain{bool}.  Primitive types are
unrestricted (non-linear).

\subsection*{structural type equivalence}
Two types are structurally equivalent if they have the same structure.
Lain uses nominal equivalence for user-defined types; the name is part of
the type identity.

\subsection*{unrestricted type}
\index{unrestricted type}
A type whose values may be used any number of times without explicit
consumption.  All primitive types are unrestricted.  Contrast with
\textit{linear type}.

% -------------------------------------------------------------------------
\section{Ownership and memory terms}
\label{sec:terms.ownership}
% -------------------------------------------------------------------------

\subsection*{active borrow}
\index{borrow!active}
A borrow that is currently in the ACTIVE phase; its subject cannot be moved
or mutated while the borrow is active (see \S\,\ref{sec:11-ownership}).

\subsection*{arena}
\index{arena}
A memory allocation strategy in which all allocations are made from a
contiguous block; the entire block is freed at once.  Lain's compiler
allocates all AST nodes from arenas.

\subsection*{borrow}
\index{borrow}
A temporary access to a value without taking ownership.  A shared borrow
grants read-only access; a mutable borrow grants read-write access
(see \S\,\ref{sec:11-ownership}).

\subsection*{borrow phase}
\index{borrow!phase}
Either RESERVED (the borrow has been registered but is not yet active,
during argument evaluation) or ACTIVE (the borrow is in force)
(see \S\,\ref{sec:11-ownership}).

\subsection*{consumed variable}
\index{variable!consumed}
A linear variable whose ownership has been transferred (via \lain{mov}).
Reading or transferring a consumed variable is a constraint violation.

\subsection*{linear state}
\index{linear state}
The tracking state of a linear variable: one of \textit{Uninit},
\textit{Unconsumed}, or \textit{Consumed} (see \S\,\ref{sec:11-ownership}).

\subsection*{move}
\index{move}
The transfer of ownership of a value from one variable or storage location
to another, performed by the \lain{mov} operator.  After a move, the
source is in the Consumed state.

\subsection*{non-lexical lifetime (NLL)}
\index{non-lexical lifetime}
A lifetime that ends at the variable's \textit{last use} rather than at the
end of its lexical scope.  Lain uses NLL to minimize the duration of borrows
(see \S\,\ref{sec:11-ownership}).

\subsection*{read-write lock invariant}
\index{read-write lock invariant}
The aliasing invariant: at any program point, for any owned value, there
exist either \textit{N}~shared borrows (N~$\geq$~0) and no mutable borrow,
or exactly one mutable borrow and no shared borrows
(see \S\,\ref{sec:11-ownership}).

\subsection*{resource}
\index{resource}
A value that represents an external entity (file handle, socket, etc.)
whose lifetime must be explicitly managed.  Resources are modelled as
linear types.

\subsection*{two-phase borrow}
\index{two-phase borrow}
A borrow that begins in the RESERVED phase during argument evaluation and
is promoted to ACTIVE after the call expression is complete
(see \S\,\ref{sec:11-ownership}).

% -------------------------------------------------------------------------
\section{Control flow and purity terms}
\label{sec:terms.purity}
% -------------------------------------------------------------------------

\subsection*{decreasing measure}
\index{decreasing measure}
An expression in a \lain{while decreasing} statement that the implementation
shall verify is non-negative when the loop condition holds and strictly
decreases on every iteration (see \S\,\ref{sec:09-statements}).

\subsection*{\texttt{func}}
\index{func}
A pure, terminating function.  A \lain{func} shall not call a \lain{proc}
or contain an unbounded \lain{while} loop (see \S\,\ref{sec:12-functions}).

\subsection*{purity}
\index{purity}
The property of a \lain{func}: its result depends only on its inputs and it
has no observable side effects on any state outside its own stack frame.

\subsection*{\texttt{proc}}
\index{proc}
A procedure that may have side effects.  A \lain{proc} may call other
\lain{proc}s, perform I/O, and contain unbounded \lain{while} loops
(see \S\,\ref{sec:12-functions}).

\subsection*{termination}
\index{termination}
The property that a function always reaches a \lain{return} statement in
a finite number of steps.  All \lain{func}s shall be verified to terminate
by the implementation (see \S\,\ref{sec:12-functions}).

% -------------------------------------------------------------------------
\section{Verification terms}
\label{sec:terms.verification}
% -------------------------------------------------------------------------

\subsection*{bounds check}
\index{bounds check}
A runtime test that verifies an array index is within the valid range
\texttt{[0, N-1]}.  A conforming implementation shall eliminate bounds
checks when VRA can statically prove the index is in range.

\subsection*{comptime}
\index{comptime}
Compile-time evaluation.  A \lain{comptime} expression or function is
evaluated by the compiler at compile time rather than at runtime
(see \S\,\ref{sec:14-generics}).

\subsection*{exhaustiveness}
\index{exhaustiveness}
The property of a \lain{case} expression in which every possible value of
the scrutinee type is covered by at least one arm
(see \S\,\ref{sec:15-match}).

\subsection*{in guard}
\index{in guard}
A use of the \lain{in} operator as a loop or conditional guard.  An in
guard of the form \lain{idx in arr} proves to the VRA that
\texttt{idx} $\in$ \texttt{[0, arr.len-1]} in the guarded body
(see \S\,\ref{sec:08-expressions}, \S\,\ref{sec:13-vra}).

\subsection*{interval}
\index{interval}
A pair $[\ell, u]$ representing the set of integers $\{x \mid \ell \leq x
\leq u\}$.  VRA represents value ranges as intervals with $-\infty$ and
$+\infty$ sentinels.

\subsection*{monomorphization}
\index{monomorphization}
The process by which a generic function is instantiated for each concrete
set of type arguments, producing a separate, non-generic translation for
each instantiation (see \S\,\ref{sec:14-generics}).

\subsection*{return constraint}
\index{return constraint}
A value range constraint on the return value of a function, declared in the
function signature as \textit{T op bound}
(see \S\,\ref{sec:13-vra}).

\subsection*{value range analysis (VRA)}
\index{value range analysis (VRA)}
A static analysis that tracks integer value ranges through the program using
interval arithmetic, used to eliminate bounds checks and prove
constraint satisfaction (see \S\,\ref{sec:13-vra}).

% -------------------------------------------------------------------------
\section{Module and name resolution terms}
\label{sec:terms.modules}
% -------------------------------------------------------------------------

\subsection*{import}
\index{import}
A directive (\lain{import module.path}) that makes the names defined in
another module available in the current module
(see \S\,\ref{sec:16-modules}).

\subsection*{module}
\index{module}
The unit of compilation in Lain.  Each source file defines one module.

\subsection*{qualified name}
\index{qualified name}
A name of the form \textit{module.identifier} that refers to a declaration
in a specific module (see \S\,\ref{sec:06-basics}).

\subsection*{Universal Function Call Syntax (UFCS)}
\index{UFCS}
\index{Universal Function Call Syntax}
A syntactic sugar by which \lain{x.f(args)} is desugared to
\lain{f(x, args)} before name resolution.  UFCS allows methods to be
defined as free functions (see \S\,\ref{sec:12-functions}).
